\documentclass[11pt]{article}

\usepackage[T1]{fontenc}
\usepackage[utf8]{inputenc}
\usepackage{mathpazo}
\usepackage[expansion=false]{microtype}
\usepackage[a4paper,margin=29mm,headheight=22pt]{geometry}
\usepackage{amsmath,amssymb,amsthm,mathtools}
\usepackage{booktabs,tabularx,array}
\usepackage[dvipsnames]{xcolor}
\usepackage{enumitem}
\usepackage{fancyhdr}
\usepackage{titlesec}
\usepackage{tcolorbox}
\usepackage{url}
\usepackage{seqsplit}
\usepackage[colorlinks=true,linkcolor=MidnightBlue,urlcolor=MidnightBlue,citecolor=MidnightBlue]{hyperref}
\usepackage[nameinlink,noabbrev]{cleveref}
\tcbuselibrary{skins,breakable}

\definecolor{navy}{HTML}{17365D}
\definecolor{bluegray}{HTML}{49677F}
\definecolor{pale}{HTML}{F3F7FA}
\definecolor{passgreen}{HTML}{1B6E43}
\definecolor{xpassblue}{HTML}{315A8A}
\definecolor{cpassviolet}{HTML}{694A86}
\definecolor{openred}{HTML}{9B2C2C}
\definecolor{rulegray}{HTML}{CBD5DF}

\hypersetup{
  pdftitle={Moving-Level Trace Towers: Barrier-Escape Addendum},
  pdfsubject={Moving-prime transpositions, density-one selection, local Faber transfer, and a large-order regular-polygon limit},
  pdfkeywords={trace polynomial, monodromy, Newton polygon, Morse polynomial, polyexponential}
}

\pagestyle{fancy}
\fancyhf{}
\fancyhead[L]{\small\color{bluegray}Moving-Level Trace Towers}
\fancyhead[R]{\small\color{bluegray}Barrier-Escape Addendum v2.4}
\fancyfoot[C]{\small\color{bluegray}\thepage}
\renewcommand{\headrulewidth}{0.3pt}
\renewcommand{\headrule}{\hbox to\headwidth{\color{rulegray}\leaders\hrule height \headrulewidth\hfill}}

\titleformat{\section}
  {\Large\bfseries\color{navy}}{\thesection.}{0.7em}{}
\titleformat{\subsection}
  {\large\bfseries\color{navy}}{\thesubsection.}{0.7em}{}
\titlespacing*{\section}{0pt}{2.4ex plus .5ex minus .2ex}{1.1ex}
\titlespacing*{\subsection}{0pt}{2.0ex plus .4ex minus .2ex}{0.8ex}

\setlength{\parindent}{0pt}
\setlength{\parskip}{0.52em}
\setlength{\emergencystretch}{2em}
\setlist[itemize]{leftmargin=1.5em,itemsep=.25em,topsep=.35em}
\setlist[enumerate]{leftmargin=1.7em,itemsep=.35em,topsep=.35em}
\renewcommand{\arraystretch}{1.18}

\newtheoremstyle{keio}
  {0.9em}{0.7em}{\itshape}{}% space, body, indent
  {\bfseries\color{navy}}{.}{0.55em}{}
\theoremstyle{keio}
\newtheorem{theorem}{Theorem}[section]
\newtheorem{corollary}[theorem]{Corollary}

\newtheoremstyle{keioinput}
  {0.9em}{0.7em}{}{}% space, body, indent
  {\bfseries\color{navy}}{.}{0.55em}{}
\theoremstyle{keioinput}
\newtheorem*{structural}{Structural input}
\newtheorem*{problem}{Window-prime selection problem}

\newcommand{\UPASS}{\textcolor{passgreen}{\textsf{U-PASS}}}
\newcommand{\XPASS}{\textcolor{xpassblue}{\textsf{X-PASS}}}
\newcommand{\CPASS}{\textcolor{cpassviolet}{\textsf{C-PASS}}}
\newcommand{\OPEN}{\textcolor{openred}{\textsf{OPEN}}}
\newcommand{\Q}{\mathbf{Q}}
\newcommand{\F}{\mathbf{F}}
\newcommand{\Gal}{\operatorname{Gal}}
\newcommand{\Ein}{\operatorname{Ein}}

\newtcolorbox{statusbox}{
  colback=pale,colframe=rulegray,boxrule=.55pt,arc=1.5mm,
  left=2.5mm,right=2.5mm,top=1.8mm,bottom=1.8mm
}
\newtcolorbox{notebox}{
  colback=white,colframe=navy,boxrule=.65pt,arc=1.5mm,
  borderline west={2.2pt}{0pt}{navy},
  left=3mm,right=2.5mm,top=2mm,bottom=2mm
}

\begin{document}

\hypersetup{pageanchor=false}
\begin{titlepage}
\thispagestyle{empty}
\vspace*{20mm}
{\color{navy}\rule{\textwidth}{1.2pt}}\\[12mm]
{\Huge\bfseries\color{navy}Moving-Level Trace Towers\par}
\vspace{4mm}
{\LARGE\bfseries Barrier-Escape Addendum\par}
\vspace{8mm}
{\large Moving-prime transpositions, density-one selection, local Faber transfer,\par
and a large-order regular-polygon limit\par}
\vfill
\begin{statusbox}
\textbf{Status legend.}
\UPASS{} denotes an unconditional theorem with a complete proof;
\XPASS{} denotes an unconditional finite theorem with exact reproducible
computation; \CPASS{} denotes a proved implication from an explicitly
displayed open hypothesis; and \OPEN{} denotes a statement not proved and
never used silently as input to an unconditional claim.
\end{statusbox}
\vfill
{\large Version 2.4\par}
{\small 1 September 2026\par}
\vspace{7mm}
{\color{navy}\rule{\textwidth}{1.2pt}}
\end{titlepage}
\hypersetup{pageanchor=true}
\pagenumbering{roman}

\begin{abstract}
Let
\[
E_q(z)=\Ein_q(z)
=\sum_{n\ge 1}\frac{(-1)^{n-1}}{n!n^q}z^n,
\qquad
P_{q,m}(X)=-m[z^m]\log(1-XE_q(z)).
\]
Earlier work proves that every $P_{q,m}$ is absolutely indecomposable and
that, for $m\ge3$, its geometric monodromy lies in $\{A_m,S_m\}$. The
unresolved odd-degree problem was to produce an odd finite branch packet
without assuming the full Trace--Morse conjecture.

This addendum develops several compatible ways around that gate. Prime
degree is Morse for every order. A strict window prime with square-free
residual derivative yields an isolated simple critical value and hence a
transposition. The criterion supplies the all-order family $m=p+2$, closes
the former first gaps $m=403$ and $m=441$, proves $S_m$ for every
$q\ge1$, $2\le m\le1001$, and proves $S_m$ for the original $q=1$ tower
through $m=10^6$. For fixed $m$, scaled traces converge effectively as
$q\to\infty$ to a binomial with regular-polygon critical data, and are
therefore eventually Morse.

The local proof depends only on a finite-jet single-spike condition and
transfers to other Faber towers; the logarithmic--Stirling tower yields an
unconditional $S_{p+2}$ family. For $q=1$, the remaining prime choice is a
shifted divisor problem for discriminants. Conditional on square-freeness
of every even residual derivative in characteristic zero, the criterion
holds for a density-one set of odd degrees, with exceptional count
$O_\varepsilon(X^{15/16+\varepsilon})$. Unconditionally, every even row has
at least half of its derivative roots simple, every prime-successor row is
fully separable, and the original tower is exactly certified through
residual row $5000$.

These results do \emph{not} prove the irrationality or transcendence of
Euler's constant. They deliberately avoid the approximation, reflexivity,
order-ledger, and limiting-measure barriers: the proofs take place at finite
level over $\Q$, finite fields, and local fields.
\end{abstract}

\tableofcontents
\newpage
\pagenumbering{arabic}

\section{Why this is the viable side of the barrier ledger}

The arithmetic-barrier analysis rules out several tempting direct attacks
on $\gamma$. The present route changes the target rather than disguising the
same obstruction.

\begin{center}
\small
\begin{tabularx}{\textwidth}{@{}>{\raggedright\arraybackslash}p{.42\textwidth}>{\raggedright\arraybackslash}X@{}}
\toprule
\textbf{Barrier in the direct special-value problem} &
\textbf{Why it is not used here}\\
\midrule
Factorial denominator cost exceeds analytic Pad\'e gain &
No rational approximation to $\gamma$ is formed.\\
Shift/reflexivity makes finite $\Ein$-data $\gamma$-blind &
The objects are finite polynomials over $\Q$, not invariant formulas intended to recover $\gamma$.\\
Adding Gamma jets creates unknown tails faster than equations &
No jet closure or zero-removal theorem is invoked.\\
Pointwise algebraic independence does not pass to a limiting measure &
No pointwise-to-limit exchange occurs.\\
Period-matrix invertibility does not imply evaluation injectivity &
No motivic or period injectivity claim occurs.\\
\bottomrule
\end{tabularx}
\end{center}

The role of the period and barrier studies is therefore architectural. They
identify finite-level monodromy as a gate where genuine progress is
possible. None of the theorems below assumes the conjectural statements of
the reference designated ``[001].''

\section{Structural input from the trace tower}

The coefficient identity
\[
[X^k]P_{q,m}(X)=\frac{m}{k}[z^m]E_q(z)^k
\]
shows that $P_{q,m}\in\Q[X]$, is monic of degree $m$, and has constant term
zero. Write
\[
G^{\mathrm{geom}}_{q,m}
=\Gal\bigl(P_{q,m}(X)-T\,/\,\overline{\Q}(T)\bigr),
\]
\[
G^{\mathrm{arith}}_{q,m}
=\Gal\bigl(P_{q,m}(X)-T\,/\,\Q(T)\bigr).
\]

We use the previously audited structural theorem.

\begin{structural}[\UPASS]
For every $q\ge1$ and $m\ge2$, the polynomial $P_{q,m}$ is
indecomposable over $\overline{\Q}$. For $m\ge3$,
\[
G^{\mathrm{geom}}_{q,m}\in\{A_m,S_m\}.
\]
For even $m$, the inertia cycle at infinity is odd, and therefore
\[
G^{\mathrm{geom}}_{q,m}=G^{\mathrm{arith}}_{q,m}=S_m.
\]
\end{structural}

The input is proved by normalized Faber right-factor rigidity, an exact
$2$-adic nonvanishing theorem for the reciprocal of $E_q$, and the
classification of primitive polynomial monodromy. The present addendum
only has to manufacture one odd finite branch cycle in odd degree.

\section{Prime degree: a complete Morse theorem}

Write
\[
A_q(z)=\frac{E_q(z)}z
=\sum_{n\ge0}a_nz^n,
\qquad
a_n=\frac{(-1)^n}{(n+1)!(n+1)^q},
\]
and
\[
P'_{q,m}(X)=\sum_{j=0}^{m-1}d_jX^j.
\]
Direct coefficient extraction gives
\begin{equation}
\boxed{d_j=m[z^{m-j-1}]A_q(z)^{j+1}.}
\label{eq:dj}
\end{equation}

\begin{theorem}[Prime-degree Morse theorem; \UPASS]\label{thm:prime}
Let $m\ge3$ be prime. For every $q\ge1$, $P_{q,m}$ is Morse and
\begin{equation}
\boxed{G^{\mathrm{geom}}_{q,m}=G^{\mathrm{arith}}_{q,m}=S_m.}
\label{eq:prime-mon}
\end{equation}
\end{theorem}

\begin{proof}
Work over $\Q_m$. Formula~\eqref{eq:dj} gives
\[
v_m(d_0)=-q,
\qquad
v_m(d_j)\ge1\quad(1\le j\le m-2),
\qquad
v_m(d_{m-1})=1.
\]
Thus the Newton polygon of $P'_{q,m}$ is the single edge
\[
(0,-q)\longrightarrow(m-1,1).
\]
Every critical point $\alpha_i$ has valuation
\[
v_m(\alpha_i)=-\frac{q+1}{m-1}.
\]
After a ramified base extension the edge residual is a binomial
$u+vY^{m-1}$. Its derivative is nonzero because $m\nmid m-1$, so its
$m-1$ roots are simple and have distinct leading residues $y_i$.

Let $c_k=[X^k]P_{q,m}$. Then $c_1=d_0$, $v_m(c_k)\ge1$ for
$2\le k\le m-1$, and $c_m=1$. At every critical point,
$\alpha_i^m$ is the unique term of smallest valuation in
$P_{q,m}(\alpha_i)$. Indeed,
\[
v_m(c_1\alpha_i)-v_m(\alpha_i^m)=1,
\]
and, for $2\le k\le m-1$,
\[
v_m(c_k\alpha_i^k)-v_m(\alpha_i^m)
\ge1+\frac{(m-k)(q+1)}{m-1}>0.
\]
The leading critical-value residues are therefore $y_i^m$. Since all
$y_i^{m-1}$ equal the same nonzero constant, $y_i^m=Cy_i$, and these
residues are pairwise distinct. All critical points are simple and all
critical values are distinct, so $P_{q,m}$ is Morse. Its finite branch
cycles are transpositions and generate $S_m$. The arithmetic group contains
the geometric one.
\end{proof}

This is an infinite, computation-free theorem: it settles every prime
degree simultaneously for every polyexponential order.

\section{A moving-prime isolated transposition}

Let $m\ge7$ be odd and choose a prime
\begin{equation}
\frac{m+1}{2}<p<m,
\qquad
r=m-p.
\label{eq:window}
\end{equation}
Then $r$ is even and $2\le r\le p-2$. The relevant valuations of the
coefficients of $A_q$ are
\[
v_p(a_n)=
\begin{cases}
0,&0\le n\le p-2,\\
-(q+1),&n=p-1,\\
-1,&p\le n\le m-1.
\end{cases}
\]
Because $m\le2p-2$, at most one exceptional factor can occur in any
composition contributing to~\eqref{eq:dj}. The exact lowest terms give the
three-edge Newton polygon
\begin{equation}
(0,-1)\longrightarrow(1,-q-1)
\longrightarrow(r,-q-1)\longrightarrow(m-1,0).
\label{eq:threeedge}
\end{equation}
The critical points split into one point of valuation $q$, $r-1$ unit
points, and $p-1$ points of valuation $-(q+1)/(p-1)$. The first point
$\alpha$ is simple, and its critical value $b=P_{q,m}(\alpha)$ satisfies
\begin{equation}
\boxed{v_p(b)=q-1.}
\label{eq:bval}
\end{equation}

Define the residual lower-row derivative
\[
R_{q,r}(X)=P'_{q,r}(X)\bmod p.
\]
Keeping the horizontal-edge principal terms gives, up to a nonzero scalar,
\begin{equation}
p^{q+1}P_{q,m}(X)\equiv\frac{X^2}{r}R_{q,r}(X)\pmod p,
\label{eq:lowerP}
\end{equation}
and differentiation gives
\begin{equation}
p^{q+1}P'_{q,m}(X)
\equiv\frac{X}{r}\bigl(2R_{q,r}(X)+XR'_{q,r}(X)\bigr)\pmod p.
\label{eq:lowerPd}
\end{equation}
The only potentially dangerous denominator when recovering $P$ from $P'$ is
removed by Frobenius:
\begin{equation}
d_{p-1}=m[z^r]A_q(z)^p
\equiv m[z^r]A_q(z^p)=0\pmod p
\qquad(0<r<p).
\label{eq:frob}
\end{equation}

\begin{theorem}[Window-prime criterion; \UPASS]\label{thm:window}
If
\begin{equation}
\gcd(P'_{q,r},P''_{q,r})=1
\quad\text{in }\F_p[X],
\label{eq:gcd}
\end{equation}
then
\begin{equation}
\boxed{G^{\mathrm{geom}}_{q,m}=G^{\mathrm{arith}}_{q,m}=S_m.}
\label{eq:window-mon}
\end{equation}
\end{theorem}

\begin{proof}
If a unit critical point $\beta$ had a critical value of elevated
valuation, \eqref{eq:lowerP}--\eqref{eq:lowerPd} would force
\[
R_{q,r}(\bar\beta)=R'_{q,r}(\bar\beta)=0,
\]
contradicting~\eqref{eq:gcd}. Thus all unit critical values have valuation
$-(q+1)$.

On the final edge the principal derivative is
$d_rX^r+mX^{m-1}$. At any of its critical points the principal critical
value is
\[
d_rX^{r+1}
\left(\frac1{r+1}-\frac1m\right)
=d_rX^{r+1}\frac{p-1}{m(r+1)},
\]
so its valuation is
\[
-\frac{m(q+1)}{p-1}<0.
\]
Hence $b$ is the unique critical value of valuation $q-1$, and it is
supported by the single simple point $\alpha$. Its inertia is a
transposition. The pre-existing indecomposability theorem gives primitivity,
and a primitive group containing a transposition is $S_m$. The arithmetic
conclusion follows.
\end{proof}

The old degree-$m$ discriminant problem has therefore descended to a gcd for
a derivative of degree $r-1=m-p-1$.

\begin{corollary}[An all-order infinite family; \UPASS]\label{cor:shifted}
For every prime $p\ge5$, every $q\ge1$, and $m=p+2$,
\begin{equation}
\boxed{G^{\mathrm{geom}}_{q,p+2}=G^{\mathrm{arith}}_{q,p+2}=S_{p+2}.}
\label{eq:shifted}
\end{equation}
\end{corollary}

Indeed $r=2$, so $P'_{q,2}\bmod p$ is linear and square-free. In
particular, the former first odd geometric gap $m=403=401+2$ and the former
first odd-square arithmetic gap $m=441=439+2$ are closed for every $q$,
without a full discriminant calculation.

\section{Exact finite rectangles}

For $p>r$, the reduction of $P_{q,r}$ depends only on $q\bmod(p-1)$. The
criterion can therefore certify infinitely many orders by a finite
residue-class computation.

\begin{theorem}[All orders through degree 1001; \XPASS]\label{thm:1001}
For every $q\ge1$ and $2\le m\le1001$,
\begin{equation}
\boxed{G^{\mathrm{geom}}_{q,m}=G^{\mathrm{arith}}_{q,m}=S_m.}
\label{eq:1001}
\end{equation}
\end{theorem}

The exact audit finds one window prime working for every residue of $q$ in
492 odd degrees. Prime-degree \cref{thm:prime} handles
$m=3,5,11,17,167,577$. Only two composite degrees require combined
certificates:
\begin{itemize}
\item At $m=27$, the primes $17,19,23$ leave the single progression
$q\equiv1075\pmod{1584}$. Modulo $28$, this progression occupies the seven
classes $3,7,11,15,19,23,27$. For each class, reduction modulo $29$ gives
two regular specializations of
$\operatorname{Disc}_X(P_{q,27}(X)-T)$ with opposite quadratic character:
\end{itemize}

\begin{center}
\small
\begin{tabular}{@{}rcc@{}}
\toprule
$q\bmod 28$ & square $(T,D(T))$ & nonsquare $(T,D(T))$\\
\midrule
3  & $(3,7)$  & $(0,17)$\\
7  & $(0,16)$ & $(4,11)$\\
11 & $(0,22)$ & $(5,27)$\\
15 & $(2,16)$ & $(1,12)$\\
19 & $(1,16)$ & $(0,18)$\\
23 & $(0,22)$ & $(1,15)$\\
27 & $(0,22)$ & $(1,3)$\\
\bottomrule
\end{tabular}
\end{center}

Hence none of the seven discriminant polynomials is a constant times a
square, supplying the required odd geometric branch packet.

\begin{itemize}
\item At $m=989$, the bad order residues for $p=919$ are odd, while those
for $p=929$ are even. Their simultaneous CRT locus is empty.
\end{itemize}

The script and captured stdout are:
\begin{itemize}
\item \nolinkurl{trace_allq_modp_audit.py}\\[-.15em]
\textcolor{bluegray}{\small SHA-256: }
\texttt{\small\seqsplit{36ba3fd004df1e5a56a7bb6c61ed49cf7542d6f7c934ba458f88c1ebe105c31c}};
\item \nolinkurl{trace_allq_modp_audit_1001.txt}\\[-.15em]
\textcolor{bluegray}{\small SHA-256: }
\texttt{\small\seqsplit{1d4867aa086b4315d45d70c52ef27090ad4d1c256cd956cb6b18d72a8be374be}}.
\end{itemize}

\begin{notebox}
This theorem asserts full symmetric monodromy. It does \emph{not} assert
that every polynomial in this rectangle is Morse.
\end{notebox}

\begin{theorem}[Original tower through degree one million; \XPASS]\label{thm:million}
For every $2\le m\le10^6$,
\begin{equation}
\boxed{G^{\mathrm{geom}}_{1,m}=G^{\mathrm{arith}}_{1,m}=S_m.}
\label{eq:million}
\end{equation}
\end{theorem}

For each odd $7\le m\le10^6$, an exact finite-field search tests the largest
strict-window prime and, only when necessary, the second largest. The largest
residual degree is $114$. The largest-prime choice fails only at
$(m,p,r)=(215,211,4)$ and $(69649,69623,26)$; the second prime works in both
cases. Even degrees are structural, and prime degrees also follow from
\cref{thm:prime} independently of the audit.

The script and captured stdout are:
\begin{itemize}
\item \nolinkurl{window_prime_selection_audit.py}\\[-.15em]
\textcolor{bluegray}{\small SHA-256: }
\texttt{\small\seqsplit{e2eb3b8a14f758d84c0d80a4ed3c0ed11a3cf7eb38e3935960be3122d8fdcbf0}};
\item \nolinkurl{window_prime_selection_audit_1000000.txt}\\[-.15em]
\textcolor{bluegray}{\small SHA-256: }
\texttt{\small\seqsplit{51d3a027e9f0ce42f4918b39058288cdfdf808e25ded994ed009a701bc7c44f5}}.
\end{itemize}

Again, this is a monodromy theorem, not an all-degree Morse theorem.

\section{Large order: complex regular polygons replace real critical points}

Put
\[
G_q(z)=\sum_{n\ge1}\frac{z^n}{n!n^q},
\qquad
h_{q;m,k}=[z^m]G_q(z)^k>0.
\]
Then
\begin{equation}
P_{q,m}(X)
=\sum_{k=1}^m(-1)^{m-k}\frac{m}{k}h_{q;m,k}X^k.
\label{eq:P-positive}
\end{equation}
If $n_1+\cdots+n_k=m$ with every $n_i\ge1$, the product of the parts
satisfies
\begin{equation}
n_1\cdots n_k\ge m-k+1.
\label{eq:parts}
\end{equation}
For $2\le r\le m-1$, define
\begin{equation}
\lambda_{m,r}=\frac{m^{(r-1)/(m-1)}}r,
\qquad
\theta_m=\max_{2\le r\le m-1}\lambda_{m,r}<1.
\label{eq:theta}
\end{equation}
The strict inequality follows from the strict decrease of
$x\mapsto\log x/(x-1)$ on $(1,\infty)$.

Let
\[
s_q=m^{-q/(m-1)},
\qquad
\widetilde P_{q,m}(y)
=m^{qm/(m-1)}P_{q,m}(s_qy).
\]
The endpoint terms $k=1,m$ remain, while
\eqref{eq:parts}--\eqref{eq:theta} suppress every intermediate term
exponentially.

\begin{theorem}[Effective large-order Morse theorem; \UPASS]\label{thm:largeq}
For every fixed $m\ge2$,
\begin{equation}
\boxed{
\widetilde P_{q,m}(y)
\longrightarrow
S_m(y):=y^m+\frac{(-1)^{m-1}}{(m-1)!}y
\quad(q\to\infty)}
\label{eq:limit}
\end{equation}
coefficientwise and locally uniformly, with error
$O_{m,R}(\theta_m^q)$ on $|y|\le R$. Consequently there is an effectively
computable $Q(m)$ such that, for every $q\ge Q(m)$,
\begin{equation}
\boxed{
P_{q,m}\text{ is Morse},
\qquad
G^{\mathrm{geom}}_{q,m}=G^{\mathrm{arith}}_{q,m}=S_m.}
\label{eq:largeq-mon}
\end{equation}
\end{theorem}

\begin{proof}
The derivative of the limit is
\[
S_m'(y)=my^{m-1}+\frac{(-1)^{m-1}}{(m-1)!}.
\]
Its roots satisfy
\[
y^{m-1}=\frac{(-1)^m}{m!},
\]
so they are the distinct vertices of a regular $(m-1)$-gon. At a critical
point $y_j$,
\[
S_m(y_j)=(-1)^{m-1}\frac{m-1}{m!}y_j,
\]
and these critical values are another set of pairwise distinct regular
polygon vertices. Thus $S_m$ is Morse. The Morse locus is Zariski open in
coefficient space; alternatively, disjoint Rouch\'e circles around these
explicit roots and their separated values give an effective bound. The
exponential coefficient estimate therefore proves~\eqref{eq:largeq-mon}.
\end{proof}

For fixed odd $m$, the limiting critical equation has no real root. Hence
sufficiently large $q$ has no real critical point at all. The smallest exact
example is $(q,m)=(2,3)$. A uniform real-transposition strategy is therefore
false; complex critical-value separation is the correct replacement.

\section{Local Faber transfer beyond the polyexponential tower}

The moving-prime mechanism is a finite-jet theorem, not a peculiarity of
the differential equation for $E_q$. Let $L$ be a complete discretely
valued characteristic-zero field with residue characteristic $p\ge5$,
uniformizer $\pi$, and residue field $k$. For
\[
f(z)=zA(z),\qquad A(z)=\sum_{n\ge0}a_nz^n,\qquad a_0=1,
\]
put $P_m^f=-m[z^m]\log(1-Xf(z))$. Suppose $m=p+r$,
$2\le r\le p-2$, and, for integers $h>s\ge0$,
\begin{equation}
\begin{array}{ll}
v(a_n)\ge0 &(0\le n\le p-2),\\
v(a_{r-1})=0,\quad v(a_{p-1})=-h,\\
v(a_n)\ge-s &(p\le n\le m-2),\\
v(a_{m-1})=-s.
\end{array}
\label{eq:spike}
\end{equation}
Let $\bar A_{<p}$ be the reduction of the integral truncation below
degree $p-1$, and define
\begin{equation}
R_r(X)=\overline{(P_r^f)'(X)}
=r\sum_{j=1}^r[z^{r-j}]\bar A_{<p}(z)^jX^{j-1}.
\label{eq:R-general}
\end{equation}

\begin{theorem}[Single-spike Faber transfer; \UPASS]\label{thm:transfer}
Under \eqref{eq:spike}, the derivative Newton polygon is exactly
\begin{equation}
(0,-s)\longrightarrow(1,-h)\longrightarrow(r,-h)
\longrightarrow(m-1,0).
\label{eq:spike-NP}
\end{equation}
If $\gcd(R_r,R_r')=1$ in $k[X]$, then $P_m^f$ has a critical value
supported by exactly one simple critical point. The three critical-point
clusters and their critical-value valuations are
\begin{equation}
\begin{array}{c|c|c}
\text{number}&v(\text{critical point})&v(\text{critical value})\\ \hline
1&h-s&h-2s\\
r-1&0&-h\\
p-1&-h/(p-1)&-mh/(p-1).
\end{array}
\label{eq:cluster-table}
\end{equation}
Thus its geometric inertia contains a transposition. If the coefficients
lie in a number field $K$, the local hypotheses hold in a completion $K_v$
or a valued extension, and $P_m^f$ is indecomposable over $\overline K$,
then both its geometric and arithmetic monodromy groups are $S_m$.
\end{theorem}

\begin{proof}
The universal identities
\begin{equation}
[X^k]P_m^f=\frac mk[z^{m-k}]A^k,
\qquad
[X^j](P_m^f)'=m[z^{m-j-1}]A^{j+1}
\label{eq:faber-universal}
\end{equation}
show that at most one copy of the unique deep coefficient $a_{p-1}$ can
enter. Exact endpoint terms give \eqref{eq:spike-NP}. If
$\eta=\overline{\pi^ha_{p-1}}$, the horizontal initial forms are
\begin{equation}
\operatorname{in}_{-h}(P_m^f)
=\frac{\bar m\eta}{r}X^2R_r(X),
\qquad
\operatorname{in}_{-h}((P_m^f)')
=\frac{\bar m\eta}{r}X(2R_r+XR_r').
\label{eq:initial-forms}
\end{equation}
The apparent coefficient $[X^p]P_m^f=d_{p-1}/p$ is integral because
$[z^r]A^p\in p\mathcal O_L$ for $0<r<p$. Consequently the gcd condition
prevents cancellation of every unit-cluster critical value. The length-one
edge gives one simple point $\alpha$, and its derivative equation yields
$v(P_m^f(\alpha))=h-2s$. The other two values in
\eqref{eq:cluster-table} follow from \eqref{eq:initial-forms} and the
final-edge binomial. Since
\[
-\frac{mh}{p-1}<-h<h-2s,
\]
the value at $\alpha$ is globally isolated. Primitive polynomial monodromy
containing a transposition is $S_m$.
\end{proof}

For $r=2$, $R_2$ is linear, so the gcd condition is automatic. This
immediately transfers the local branch construction, for integers
$\tau\ge1$ and $\sigma,\tau\ge1$, respectively, to
\begin{equation}
\operatorname{Li}_\tau(z)=\sum_{N\ge1}\frac{z^N}{N^\tau}
\quad\text{and}\quad
f_{\sigma,\tau}(z)=\sum_{N\ge1}\frac{z^N}{(N!)^\sigma N^\tau},
\label{eq:transfer-families}
\end{equation}
with depths $(h,s)=(\tau,0)$ and $(\sigma+\tau,\sigma)$ over $\Q_p$,
respectively. These give a transposition; full $S_m$ additionally requires
indecomposability.

There is one unconditional new global example. For
$f(z)=-\log(1-z)$,
\begin{equation}
P_m^f(X)=\sum_{k=1}^m
\frac{(k-1)!\,|s(m,k)|}{(m-1)!}X^k.
\label{eq:stirling}
\end{equation}
Moreover
\[
\frac z{-\log(1-z)}=\int_0^1(1-z)^t\,dt
\]
has a nonzero coefficient in every positive degree. Hence every Faber
right-factor obstruction $d_b=[z^b](1/f)=[z^{b+1}](z/f)$ is nonzero, so
all the polynomials in \eqref{eq:stirling} are absolutely indecomposable.
\Cref{thm:transfer} therefore gives
\begin{equation}
\boxed{
\Gal(P_{p+2}^f(X)-T/\overline{\Q}(T))
=\Gal(P_{p+2}^f(X)-T/\Q(T))=S_{p+2}
\quad(p\ge5\text{ prime})}
\label{eq:stirling-mon}
\end{equation}
for the logarithmic--Stirling trace tower. No literature-priority claim is
made without a dedicated search.

\section{Shifted discriminants, partial separability, and conditional density one}

For the original order $q=1$, put
\begin{equation}
F_r(X)=(r!)^2P'_{1,r}(X)\in\mathbf Z[X],
\qquad
\Delta_r=\operatorname{Disc}(F_r).
\label{eq:Delta}
\end{equation}
The integrality follows termwise from $\prod n_i!\mid r!$ and
$\prod n_i\mid r!$ for every composition $n_1+\cdots+n_k=r$. If $p>r$,
the leading coefficient $r(r!)^2$ is a unit modulo $p$, so
\begin{equation}
\boxed{
P'_{1,r}\bmod p\text{ is not square-free}
\quad\Longleftrightarrow\quad p\mid\Delta_r.}
\label{eq:Delta-equivalence}
\end{equation}
Coefficient majorization by $(e^z-1)^k$, followed by Mahler's
discriminant bound, gives whenever $\Delta_r\ne0$
\begin{equation}
\log|\Delta_r|=O(r^2\log r),
\qquad
\omega(\Delta_r)=O(r^2\log r).
\label{eq:Delta-bound}
\end{equation}

\begin{notebox}
\textbf{Hypothesis DS (\OPEN).} Every even residual derivative
$P'_{1,r}$, $r\ge2$, is square-free over $\Q$, equivalently
$\Delta_r\ne0$.
\end{notebox}

\begin{theorem}[Unconditional partial separability; \UPASS]\label{thm:partial-sep}
Let $r=2s$. For every $q\ge1$, the derivative $P'_{q,r}$ has at least
$s$ simple roots. After the dyadic normalization
\[
c_q=2^{q+1},\qquad
\widehat P_{q,r}(X)=c_q^rP_{q,r}(X/c_q),
\]
these are exactly the unit roots, while every possible repeated root lies
in the strict positive-valuation disk. Moreover, for every odd prime $p$
and every $q\ge1$,
\begin{equation}
\boxed{P'_{q,p+1}(X)\text{ is square-free over }\Q.}
\label{eq:prime-successor}
\end{equation}
\end{theorem}

\begin{proof}
Put $f_q(z)=E_q(c_qz)/c_q$. Exact coefficient valuations give the
order-independent reduction
\[
f_q(z)\equiv z+z^2\pmod2.
\]
If $a=v_2(r)$, coefficient extraction then yields
\begin{equation}
2^{-a}\widehat P'_{q,2s}(X)\equiv
X^{s-1}\Lambda_{2s+1}(X)\pmod2,
\label{eq:lambda-reduction}
\end{equation}
\begin{equation}
\Lambda_{2s+1}(X)=
\sum_{j=0}^s\binom{s+j}{s-j}X^j.
\label{eq:lambda}
\end{equation}
The Lucas recurrence $L_0=0,L_1=X$,
$L_n=X(L_{n-1}+L_{n-2})$ gives
$L_{2s+1}=X^{s+1}\Lambda_{2s+1}$. Its nonzero roots are
$X=\zeta+\zeta^{-1}$ for nontrivial $(2s+1)$-st roots of unity, modulo
$\zeta\leftrightarrow\zeta^{-1}$. Hence $\Lambda_{2s+1}$ has $s$
distinct nonzero roots and simple Hensel lifts, proving the first assertion.

For $m=p+1$, the $p$-adic Newton polygon of $P'_{q,m}$ is
\begin{equation}
(0,-1)\longrightarrow(1,-q-1)\longrightarrow(p,0).
\label{eq:prime-successor-NP}
\end{equation}
The first edge is linear; after adjoining a uniformizer for the possibly
fractional second slope, its nonzero-root initial binomial has exponent gap
$p-1$, prime to $p$. Thus all roots are simple.
\end{proof}

\begin{corollary}[Exact residual rectangle; \XPASS]\label{cor:residual-5000}
For the original order, exact finite-field computation proves
\begin{equation}
\boxed{P'_{1,r}\text{ is square-free for every even }2\le r\le5000.}
\label{eq:residual-5000}
\end{equation}
\end{corollary}

At the single good prime $998244353$, exact convolution and polynomial gcd
arithmetic certify \eqref{eq:residual-5000}. Leading coefficients stay
nonzero modulo this prime; for odd $3\le s\le4999$, those of the polar and
its derivative are $s+1$ and $(s-1)(s+1)$, respectively. Thus the relevant
degrees are preserved. The same run, together with the trivial constant
case $s=1$, also certifies every odd polar $2H_s+XH_s'$ for
$1\le s\le4999$, where
$H_s=\sum_k[z^s](-E_1(-z))^kX^{k-1}$.

The exact audit has SHA-256 hashes
\begin{itemize}
\item \nolinkurl{even_derivative_separability_audit.cpp}\\[-.15em]
\textcolor{bluegray}{\small SHA-256: }
\texttt{\small\seqsplit{d9c6f7c6438a9e7bf5d0aed50bbce727f460efcacc720f520c97b444b3f83b4d}};
\item \nolinkurl{even_derivative_separability_audit_5000.txt}\\[-.15em]
\textcolor{bluegray}{\small SHA-256: }
\texttt{\small\seqsplit{1e4b4073527d7d2543c8a087c37f0888fdbd99c7ca2b7880a418e7c533e6617b}}.
\end{itemize}
This does not prove DS: in row $r=10$, a later inner Newton face already
has inseparable residual polynomial $1+Y^2$ over $\F_2$.

\begin{theorem}[Conditional density-one selection; \CPASS]\label{thm:density-one}
Assume DS. Let $\mathcal E(X)$ be the odd $m\le X$ for which no prime
$(m+1)/2<p<m$ passes the residual gcd criterion \eqref{eq:gcd}. Then, for
every $\varepsilon>0$,
\begin{equation}
\boxed{\#\mathcal E(X)\ll_\varepsilon X^{15/16+\varepsilon}.}
\label{eq:density-one}
\end{equation}
In particular, under DS the original trace tower has symmetric geometric
and arithmetic monodromy for a density-one set of odd degrees.
\end{theorem}

\begin{proof}
In a dyadic interval $[X,2X]$, let $p<m$ be the largest prime and write
$r=m-p$. Set $H=X^{5/16}$. If $r\le H$ and $m$ is exceptional, then
\eqref{eq:Delta-equivalence}--\eqref{eq:Delta-bound} give
\[
\#\{m:r\le H,\ p\mid\Delta_r\}
\ll\sum_{r\le H}r^2\log r
\ll H^3\log H.
\]
The remaining integers lie more than $H$ places into a prime gap. The
\href{https://arxiv.org/abs/1201.1787}{Maynard--Peck mean-square gap bound}
\[
\sum_{p_n\le x}(p_{n+1}-p_n)^2
\ll_\varepsilon x^{5/4+\varepsilon}
\]
implies that their number is
$O_\varepsilon(X^{5/4+\varepsilon}/H)$. Balancing the two estimates gives
\eqref{eq:density-one}; dyadic summation completes the proof.
\end{proof}

The exact equivalence \eqref{eq:Delta-equivalence} also explains why
prime-gap bounds alone do not prove selection in every degree: one must
exclude the correlated divisibilities $m-r\mid\Delta_r$ for several
offsets. Already at $m=215$, the largest window prime $211$ is bad in row
$r=4$, although the next prime works. Thus a canonical nearest-prime proof
is false even though the selectable-prime criterion survives.

\section{Exact negative results that prevent a false proof}

The failed routes are useful because they make the remaining gap precise.
\begin{enumerate}
\item A simple real critical point does not suffice: another critical point
can share its critical value. For $q=1$, every odd degree has a simple
positive critical point, but global value separation remains necessary.
\item The natural coefficient triangle is not totally positive. At $q=1$
it already has an exact negative $9\times9$ minor. A proof cannot cite global
$TP_\infty$.
\item The dyadic derivative analysis isolates a simple outer critical
cluster, but coefficient Newton polygons alone do not exclude an exact
collision with a shallower cluster.
\item The moving-prime criterion is sufficient, not automatic. For odd
composite $m$, Bertrand applied to $(m+1)/2$ supplies a prime
$(m+1)/2<p<m$; compositeness excludes the only possible upper endpoint
$p=m$. Bertrand says nothing, however, about residual square-freeness for
an eligible prime.
\end{enumerate}

\section{The remaining theorem gate}

The cleanest unresolved statement is now finite-field rather than
transcendental.

\begin{problem}[\OPEN]
Prove that for every $q\ge1$ and every sufficiently large odd composite $m$,
at least one prime $p=p(q,m)$
\[
\frac{m+1}{2}<p<m
\]
makes
\[
\gcd(P'_{q,m-p},P''_{q,m-p})=1\pmod p.
\]
\end{problem}

For the original order $q=1$, exact computation verifies this through
$m=10^6$, always using one of the two largest eligible primes. Under DS,
\cref{thm:density-one} proves it for a density-one set, but no argument
currently promotes that finite fact to all $m$. A proof would close the
remaining odd composite degree problem by \cref{thm:window}.

An alternative is to continue the polar expansion at a multiple residual
root and prove that its lifted critical value cannot meet the isolated value
of~\eqref{eq:bval}. This is the precise local collision problem left by the
present work.

\section{Integration map for the main manuscript}

This addendum is designed to be inserted immediately after the existing
section on universal level-trace polynomials. A conservative revision should
make the following changes elsewhere in the manuscript.
\begin{enumerate}\sloppy
\item Add to the abstract the prime-degree Morse theorem, the moving-prime
criterion, the all-order rectangle $m\le1001$, and the original-order range
$m\le10^6$.
\item Replace every statement that $m=403$ is the first unresolved geometric
degree and $m=441$ the first unresolved arithmetic odd-square degree.
\Cref{cor:shifted} closes both for every $q$.
\item Retain the existing arithmetic-barrier section, but add a forward
reference explaining that the monodromy work escapes the barrier by staying
at finite algebraic level. Do not present it as a route to the transcendence
of $\gamma$.
\item Keep the distinction between the Morse theorem and the group theorem:
prime degrees and sufficiently large orders are Morse; the finite rectangles
assert $S_m$ and need not be labelled Morse.
\item Replace the old full-discriminant computational emphasis by the
residual derivative certificate. The full discriminant is retained only for
the exceptional $m=27$ residue progression.
\item In the final open-problem list, promote the window-prime selection
problem and the lifted residual-collision problem. Keep $\gamma$'s arithmetic
nature explicitly open.
\item Add the finite-jet Faber transfer theorem as the conceptual endpoint
of the local argument, with the logarithmic--Stirling family as an
unconditional application.
\item State the shifted-discriminant formulation and the DS-conditional
density-one theorem in a separate paragraph whose conditional status is
visible in both the theorem title and claim table.
\item {\raggedright Record the unconditional dyadic half-separation theorem,
the all-order prime-successor residual family, and the exact DS audit through
5000; retain full DS as an open hypothesis.\par}
\end{enumerate}

\section{Claim ledger and realistic value}

\begin{center}
\small
\begin{tabularx}{\textwidth}{@{}>{\raggedright\arraybackslash}X>{\centering\arraybackslash}p{.16\textwidth}>{\raggedright\arraybackslash}p{.27\textwidth}@{}}
\toprule
\textbf{Result} & \textbf{Status} & \textbf{Scope}\\
\midrule
All-$(q,m)$ indecomposability and $A_m/S_m$ reduction & \UPASS & $q\ge1,m\ge2$; reduction for $m\ge3$\\
Prime-degree Morse and $S_m$ & \UPASS & Every $q$, every prime $m\ge3$\\
Window-prime isolated-transposition criterion & \UPASS & Every qualifying $(q,m,p)$\\
Shifted-prime family $m=p+2$ & \UPASS & Every $q$, infinitely many odd $m$\\
Large-order Morse theorem & \UPASS & Fixed $m$, all sufficiently large $q$\\
Single-spike Faber transfer theorem & \UPASS & Transposition under \eqref{eq:spike}; $S_m$ with indecomposability\\
Logarithmic--Stirling $S_{p+2}$ family & \UPASS & Every prime $p\ge5$\\
All orders through $m=1001$ & \XPASS & Every $q\ge1$\\
Original tower through $m=10^6$ & \XPASS & $q=1$\\
At least half of every even derivative row simple & \UPASS & Every $q\ge1$, every even row\\
Prime-successor residual rows square-free & \UPASS & $r=p+1$, every odd prime $p$, every $q$\\
Residual DS through row 5000 & \XPASS & $q=1$, every even $r\le5000$\\
Density-one odd-degree selection & \CPASS & $q=1$, conditional on DS\\
Characteristic-zero residual separability (DS) & \OPEN & Every even residual row\\
All odd composite degrees & \OPEN & Selection/local-collision lemma missing\\
Irrationality or transcendence of $\gamma$ & \OPEN & Untouched\\
\bottomrule
\end{tabularx}
\end{center}

The package is a material increase in the paper's value: it replaces a
global discriminant search by a local recursive certificate, proves new
infinite regimes in two trace families, transfers the mechanism beyond
$\Ein_q$, and expands exact coverage by several orders of magnitude. The
shifted-discriminant formulation also reduces a large part of the asymptotic
obstruction to the explicit hypothesis DS. It is not yet an Annals-level
closure of the full two-parameter problem. The natural high-tier target is
DS together with window-prime selection, or an unconditional global
critical-value separation theorem; either would turn the present
architecture into a full all-odd result.

\section{Reproducibility map}

\begin{itemize}\raggedright
\item \nolinkurl{TRACE_MONODROMY_ALL_Q.md}: structural Faber and monodromy input.
\item \nolinkurl{ODD_TRACE_MODP_ROUTE.md}: prime-degree and moving-prime proofs.
\item \nolinkurl{ODD_TRACE_PADIC_ROUTE.md}: independent derivation of the
isolated-branch criterion and the $q=1$ audit.
\item \nolinkurl{ODD_TRACE_REAL_CRITICAL_ROUTE.md}: large-order
regular-polygon theorem and exact failure of the uniform real route.
\item \nolinkurl{TRACE_DERIVATIVE_STRUCTURE_ROUTE.md}: second-kind Faber and
dyadic local critical-cluster identities.
\item \nolinkurl{ABSTRACT_FABER_MOVING_PRIME_TRANSFER.md}: complete local
transfer theorem and polylogarithmic, hypergeometric, and logarithmic
applications.
\item \nolinkurl{ABSTRACT_FABER_TRANSFER_REFEREE.md}: independent hostile
audit of the transfer theorem.
\item \nolinkurl{WINDOW_PRIME_SELECTION_AUDIT.md}: shifted-discriminant
reduction, conditional density-one theorem, and exact million-degree audit.
\item \nolinkurl{WINDOW_PRIME_SELECTION_REFEREE.md}: independent audit of the
selection theorem and computation.
\item \nolinkurl{EVEN_DERIVATIVE_SEPARABILITY.md}: dyadic half-separation,
prime-shift families, and exact residual-row audit through 5000.
\item \nolinkurl{EVEN_DERIVATIVE_SEPARABILITY_REFEREE.md}: independent audit
of the DS partial results and finite computation.
\item \nolinkurl{trace_allq_modp_audit.py}: all-order finite-field audit
through degree 1001.
\item \nolinkurl{window_prime_selection_audit.py}: original-order selection
audit through degree one million.
\end{itemize}

\vfill
\begin{statusbox}
\small
\textbf{Status reminder.} \UPASS{} means unconditional theorem with complete
proof; \XPASS{} means unconditional finite theorem with exact reproducible
computation; \CPASS{} means a proved implication from an explicitly stated
open hypothesis; \OPEN{} means not proved and never silently used as theorem
input.
\end{statusbox}

\end{document}
